\chapter{Result: Theorem 1.6 (Rectangle Partition Bound)}

\begin{definition}[Leaf rectangles]
  \label{def:leaf_rectangles}
  \lean{DetProtocol.leafRectangles}
  \leanok
  \uses{def:det_protocol, def:rectangle}
  The leaf rectangles of a protocol are the input sets reaching each leaf.
\end{definition}

\begin{lemma}[Leaf rectangles form a rectangle partition]
  \label{lem:leaf_partition}
  \lean{DetProtocol.leafRectangles_cover}
  \lean{DetProtocol.leafRectangles_disjoint}
  \lean{DetProtocol.leafRectangles_isRectangle}
  \leanok
  \uses{def:leaf_rectangles, def:mono_rect_partition}
  Leaf rectangles cover $X \times Y$, are pairwise disjoint, and are rectangles.
\end{lemma}

\begin{lemma}[Leaf count bound]
  \label{lem:leaf_card}
  \lean{DetProtocol.leafRectangles_card}
  \leanok
  \uses{def:leaf_rectangles, def:protocol_run}
  A protocol of complexity $c$ has at most $2^c$ leaf rectangles.
\end{lemma}

\begin{theorem}[Theorem 1.6]
  \label{thm:rect_partition}
  \lean{DetProtocol.rectangle_partition}
  \lean{DetProtocol.mono_rectangle_partition_of_det_cc}
  \leanok
  \uses{lem:leaf_partition, lem:leaf_card, def:mono_rect_partition}
  If a protocol computes $g$ with complexity $c$, then $X \times Y$ admits a
  monochromatic rectangle partition for $g$ of size at most $2^c$.
\end{theorem}
